\documentclass[12pt]{article}
\usepackage[T1]{fontenc}
\usepackage[utf8]{inputenc}
\usepackage{lmodern}
\usepackage{textcomp}
\usepackage{enumitem}
\usepackage{geometry}
\usepackage{graphicx}
\usepackage{amsmath, amssymb}
\geometry{margin=1in}
\begin{document}
\begin{center}
Biology - Undergraduate Level Assessment 1 \
Total Marks: 40 \
Answer all questions.
\end{center}

\section*{Multiple Choice (10 marks)}
\begin{enumerate}[label=\arabic*.]
\item Which of the following characteristics is predicted for an early-successional plant community?
\begin{enumerate}[label=(\alph*)]
    \item High detrital biomass
    \item High number of predatory species
    \item High presence of fully matured plant species
    \item High frequency of R-selected species
    \item High rates of soil nutrient depletion
\end{enumerate}
\item Mitochondria isolated and placed in a buffered solution with a low pH begin to manufacture ATP. Which of the following is the best explanation for the effect of low external pH?
\begin{enumerate}[label=(\alph*)]
    \item It increases the acid concentration in the mitochondria matrix.
    \item It decreases the acid concentration in the mitochondria matrix.
    \item It increases the concentration of H+ in the mitochondria matrix.
    \item It decreases the OH- concentration in the mitochondria matrix.
    \item It increases diffusion of H+ from the intermembrane space to the matrix.
\end{enumerate}
\item In reference to a segment of DNA, which of the following molecules contains the fewest number of nucleotides?
\begin{enumerate}[label=(\alph*)]
    \item the tRNA transcript from the original DNA
    \item a single strand of the original DNA segment after a substitution mutation
    \item the final processed mRNA made from the original DNA
    \item the primary RNA transcript (after splicing) from the original DNA
    \item a single strand of the original DNA segment after a duplication mutation
\end{enumerate}
\item All of the following might be found in connective tissue EXCEPT
\begin{enumerate}[label=(\alph*)]
    \item thrombin
    \item glycosaminoglycans
    \item collagens
    \item fibroblasts
\end{enumerate}
\item What properties of water make it an essential component of living matter?
\begin{enumerate}[label=(\alph*)]
    \item Water is useful just because of its cooling properties.
    \item Water is only important for hydration.
    \item Water's role is limited to acting as a chemical reactant in cellular processes.
    \item Water is important primarily because it contributes to the structural rigidity of cells.
    \item Water's only function is to provide a medium for aquatic organisms.
\end{enumerate}
\end{enumerate}

\section*{True / False (10 marks)}
\textit{Answer True or False}
\begin{enumerate}[label=\arabic*.]
\setcounter{enumi}{5}
\item True or False: The neurilemma sheath provides a channel for the axon to grow back to its former position during nerve regeneration.
\item True or False: Genetic engineering is the main source of variation observed among different human populations.
\item True or False: Removing the macronucleus from a paramecium causes cell death due to the loss of essential functions.
\item True or False: Water moves from the surrounding pure water into the potato, causing the potato to swell.
\item True or False: An increase in the amount of CO 2 in the blood does not influence the regulation of respiration.
\end{enumerate}

\section*{Long Form Response (20 marks)}
\textit{Answer each question with a clear and complete explanation.}
\begin{enumerate}[label=\arabic*.]
\setcounter{enumi}{10}
\item Discuss the role of hydrochloric acid in the stomach and analyze how its absence due to abnormal conditions could impact the digestive process and overall gut health.
\item Discuss the structure and function of the parathyroid glands in humans, highlighting their role in calcium regulation and the implications of their dysfunction on human health.
\end{enumerate}

\vfill
\noindent\textit{End of Paper}
\end{document}